\documentclass[11pt,a4paper]{article}
\usepackage{fontspec}
\usepackage{polyglossia}
\setmainlanguage{portuguese}
\usepackage[margin=2.5cm]{geometry}
\usepackage{amsmath}
\usepackage{booktabs}
\usepackage{tabularx}
\usepackage{listings}
\usepackage{xcolor}

\lstset{
  language=C++,
  basicstyle=\ttfamily\small,
  keywordstyle=\bfseries,
  commentstyle=\color{gray}\itshape,
  numbers=left, numberstyle=\tiny\color{gray},
  frame=single, breaklines=true, columns=fullflexible,
  literate={º}{{$^\circ$}}1
}

\title{Trabalho Final -- Inclinômetro com ADXL335}
\author{Bernardo}
\date{\today}

\begin{document}
\maketitle

\section{Objetivo}
Medir o ângulo de inclinação em torno do eixo X com o acelerômetro
analógico ADXL335 (módulo GY-61), mostrar o valor em graus num OLED
SSD1306 e acender um LED quando a inclinação chega a 30$^\circ$ para
qualquer lado.

\section{Critérios do projeto}
A Tabela~\ref{tab:criterios} lista o que o sistema precisa cumprir e
onde cada ponto é resolvido. As seções seguintes explicam cada um.

\begin{table}[h]
\centering
\begin{tabularx}{\textwidth}{>{\raggedright}p{4.2cm}X}
\toprule
Critério & Como é atendido \\
\midrule
Medir o ângulo em X & $\operatorname{atan2}(a_y, a_z)$ com os eixos Y e Z
  (Seção~\ref{sec:angulo}) \\
Exatidão independente do USB & leitura dos 3,3~V em A3 e divisão
  raciométrica (Seção~\ref{sec:ref}) \\
Leitura sem ruído & média de 16 amostras + filtro passa-baixa
  (Seções~\ref{sec:media} e \ref{sec:filtro}) \\
Número estável na tela & histerese de 0,8$^\circ$ (Seção~\ref{sec:hist}) \\
Alarme a 30$^\circ$ & LED em D8 aceso se $|\theta| \ge 30^\circ$ \\
Segurança elétrica do OLED & pull-ups internos de 5~V desligados
  (Seção~\ref{sec:hw}) \\
Ajuste a cada sensor & constantes \texttt{ZERO} e \texttt{ESCALA}
  (Seção~\ref{sec:calib}) \\
\bottomrule
\end{tabularx}
\caption{Critérios e solução adotada.}
\label{tab:criterios}
\end{table}

\section{Material e ligações}
Arduino Uno (ATmega328P, 16~MHz), módulo GY-61 (ADXL335), OLED
128$\times$64 I\textsuperscript{2}C (endereço 0x3C), LED e resistor de
220~$\Omega$.

\begin{center}
\begin{tabular}{ll}
\toprule
Pino do Uno & Ligação \\
\midrule
3,3 V / GND & VCC e GND do GY-61 e do OLED \\
A1, A2 & saídas Y e Z do GY-61 \\
A3 & trilho de 3,3 V (referência) \\
A4 / A5 & SDA / SCL do OLED \\
D8 & 220~$\Omega$ $\rightarrow$ LED $\rightarrow$ GND \\
\bottomrule
\end{tabular}
\end{center}

A saída X do GY-61 não é ligada: para inclinar em torno do eixo X só
importam Y e Z.

\section{Fundamentos}

\subsection{Como o sensor responde}
O ADXL335 mede aceleração em três eixos e entrega uma tensão para cada
um. Parado, a única aceleração que ele sente é a gravidade (1~g). Com
o eixo na horizontal (0~g nele), a saída fica na metade da alimentação:
1,65~V com 3,3~V. Cada 1~g desloca a saída cerca de um décimo da
alimentação, 0,33~V. Como esses valores acompanham a tensão de
alimentação, o sensor é chamado de \emph{raciométrico}.

\subsection{O ADC do Uno}
O conversor analógico-digital (ADC) de 10 bits transforma a tensão do
pino num número $N$ de 0 a 1023:
\[
  N = \left\lfloor \frac{V_\text{pino}}{V_\text{ref}} \times 1024 \right\rfloor
\]
onde $V_\text{ref}$ é a referência do ADC, por padrão os 5~V do USB.

\subsection{Por que ler os 3,3~V em A3}\label{sec:ref}
Os 5~V do USB não são exatos: variam de 4,7 a 5,1~V conforme a porta e
o cabo. Se o programa usasse direto a leitura do eixo, a mesma tensão
viraria números diferentes. Exemplo: 0,1~V a menos no USB (2\%) muda a
leitura do eixo Z em cerca de 10 contagens, o que equivale a 0,1~g e
dá perto de 6$^\circ$ de erro.

A solução é medir, com o mesmo ADC, a própria alimentação do sensor
(3,3~V em A3) e dividir uma leitura pela outra:
\[
  \frac{N_\text{eixo}}{N_\text{A3}}
  \approx \frac{V_\text{eixo}/V_\text{ref}}{3{,}3/V_\text{ref}}
  = \frac{V_\text{eixo}}{3{,}3}
\]
O $V_\text{ref}$ some da conta. Multiplicando por 1024, o resultado
fica na mesma faixa de uma leitura comum, mas agora relativo aos
3,3~V. Como o sensor também é relativo aos 3,3~V, os valores teóricos
ficam fixos: 0~g vale $1024/2 = 512$ e 1~g vale $1024/10 = 102{,}4$.

\subsection{Cálculo do ângulo}\label{sec:angulo}
Girando o sensor em torno do eixo X por um ângulo $\theta$, a gravidade
se divide entre Y e Z:
\[
  a_y = \sin\theta \cdot 1\,\text{g}, \qquad a_z = \cos\theta \cdot 1\,\text{g}
\]
Logo $\theta = \arctan(a_y / a_z)$. O programa usa
$\operatorname{atan2}(a_y, a_z)$, que faz a mesma conta mas olha o
sinal dos dois argumentos e devolve o ângulo no quadrante certo, de
$-180^\circ$ a $+180^\circ$. Ela também não divide por zero quando
$a_z = 0$ (sensor a 90$^\circ$).

Uma vantagem: como entra a \emph{razão} entre $a_y$ e $a_z$, um erro de
escala igual nos dois eixos se cancela. Por isso o ângulo é bem menos
sensível à constante \texttt{ESCALA} do que ao \texttt{ZERO}.

\textbf{Exemplo.} Inclinado 30$^\circ$: $a_y = 0{,}5$~g e
$a_z = 0{,}866$~g. As leituras corrigidas seriam
$512 + 0{,}5 \times 102{,}4 = 563{,}2$ e
$512 + 0{,}866 \times 102{,}4 = 600{,}7$.
Voltando: $\operatorname{atan2}(0{,}5;\ 0{,}866) = 30^\circ$.

\subsection{Média de 16 leituras}\label{sec:media}
Cada leitura do ADC tem ruído de uma ou duas contagens. Somar 16
leituras e dividir por 16 reduz esse ruído cerca de 4 vezes
($\sqrt{16}$). A soma máxima é $16 \times 1023 = 16\,368$, que cabe
num \texttt{unsigned int} de 16 bits (até 65\,535).

Antes das 16 leituras há uma leitura descartada. O ADC tem um só
conversor e um seletor (multiplexador) que escolhe o pino; logo após
trocar de pino, a primeira conversão ainda pode trazer resto do pino
anterior.

\subsection{Filtro passa-baixa}\label{sec:filtro}
Mesmo com a média, o ângulo ainda oscila. O filtro faz o valor
guardado andar só uma fração $\alpha = 0{,}05$ (5\%) em direção à
leitura nova:
\[
  \theta_\text{filtrado} \leftarrow \theta_\text{filtrado}
    + \alpha\,(\theta_\text{novo} - \theta_\text{filtrado})
\]
É uma média ponderada: 95\% do valor antigo e 5\% do novo. Ruído
rápido, que sobe e desce, quase não move o resultado; uma mudança real
e duradoura chega aos poucos. Após $n$ ciclos resta $0{,}95^n$ da
diferença inicial: 37\% após 20 ciclos e 5\% após 60 ciclos.

O custo é atraso. Cada ciclo dura os 50~ms do \texttt{delay} mais o
tempo de ler o ADC e de enviar a tela (estimado em algumas dezenas de
milissegundos, não medido). Assim, a tela leva da ordem de 1 a 2~s
para acompanhar a maior parte de um movimento e alguns segundos para
assentar de vez.

\subsection{Histerese}\label{sec:hist}
A tela mostra um número inteiro. Se o ângulo filtrado estivesse em
20,5$^\circ$, um arredondamento simples ficaria pulando entre 20 e 21.
A histerese evita isso: o número mostrado só é atualizado quando o
ângulo filtrado se afasta mais de 0,8$^\circ$ dele. Exemplo: mostrando
20, o valor só muda se o ângulo passar de 20,8 ou cair abaixo de 19,2.
Como 0,8 é maior que 0,5 (meio grau do arredondamento), sobra uma
margem em que pequenas oscilações não mudam nada.

O LED usa o número mostrado, então ele também herda essa estabilidade
e não pisca perto dos 30$^\circ$.

\section{O código, parte por parte}

\subsection{Bibliotecas e constantes}
\begin{lstlisting}
#include <Wire.h>
#include <Adafruit_GFX.h>
#include <Adafruit_SSD1306.h>

const uint8_t PINO_Y = A1, PINO_Z = A2, PINO_REF = A3, LED = 8;
const float ZERO = 512;      // leitura em 0 g
const float ESCALA = 102.4;  // leitura por 1 g
const int LIMIAR = 30;       // graus
const float ALFA = 0.05;     // peso da leitura nova no filtro
\end{lstlisting}
\texttt{Wire} é a biblioteca do I\textsuperscript{2}C (o barramento de
dois fios do OLED). \texttt{Adafruit\_GFX} desenha texto e formas e
\texttt{Adafruit\_SSD1306} fala com o controlador da tela. As
constantes reúnem num só lugar tudo o que se pode querer ajustar:
pinos, calibração, limiar do LED e força do filtro.

\subsection{Variáveis globais}
\begin{lstlisting}
Adafruit_SSD1306 tela(128, 64, &Wire, -1);
float angulo = 0;
int mostrado = 0;
\end{lstlisting}
\texttt{tela} é o objeto do OLED: 128$\times$64 pixels, no barramento
\texttt{Wire}, sem pino de reset ($-1$). \texttt{angulo} guarda o
valor filtrado e \texttt{mostrado} o inteiro da tela. Elas ficam fora
das funções porque precisam manter o valor de um ciclo do
\texttt{loop} para o próximo: o filtro e a histerese dependem do valor
anterior.

\subsection{Função \texttt{le}: média de 16 leituras}
\begin{lstlisting}
float le(uint8_t pino) {
  analogRead(pino);                 // descarta a primeira
  unsigned int soma = 0;
  for (uint8_t i = 0; i < 16; i++) soma += analogRead(pino);
  return soma / 16.0;
}
\end{lstlisting}
A primeira \texttt{analogRead} só troca o canal e é jogada fora
(Seção~\ref{sec:media}). Dividir por \texttt{16.0}, e não por
\texttt{16}, força a conta em ponto flutuante e preserva as casas
decimais da média.

\subsection{Função \texttt{leG}: leitura em g}
\begin{lstlisting}
float leG(uint8_t pino, float ref) {
  return (le(pino) / ref * 1024 - ZERO) / ESCALA;
}
\end{lstlisting}
É a fórmula da Seção~\ref{sec:ref} em uma linha, em três passos:
\begin{enumerate}
  \item $N/N_\text{A3} \times 1024$: leitura relativa
    aos 3,3~V, livre do erro do USB;
  \item $-\,$\texttt{ZERO}: tira o ponto de 0~g, deixando só o desvio;
  \item $/\,$\texttt{ESCALA}: converte contagens em g.
\end{enumerate}
O \texttt{ref} vem de fora porque é lido uma vez por ciclo e usado nos
dois eixos.

\subsection{\texttt{setup}: executado uma vez}
\begin{lstlisting}
void setup() {
  pinMode(LED, OUTPUT);
  tela.begin(SSD1306_SWITCHCAPVCC, 0x3C);
  digitalWrite(SDA, LOW);   // desliga pull-up interno
  digitalWrite(SCL, LOW);
  tela.cp437(true);
  tela.setTextColor(SSD1306_WHITE);
  tela.setTextSize(2);
}
\end{lstlisting}
\begin{itemize}
  \item \texttt{pinMode} configura D8 como saída para o LED.
  \item \texttt{tela.begin} inicia o I\textsuperscript{2}C e a tela no
    endereço 0x3C. \texttt{SWITCHCAPVCC} diz que a tensão alta do
    painel é gerada dentro do próprio módulo.
  \item Os dois \texttt{digitalWrite(\dots, LOW)} desligam os pull-ups
    internos (Seção~\ref{sec:hw}). Num pino configurado como entrada,
    escrever LOW desliga o pull-up em vez de mandar 0~V.
  \item \texttt{cp437(true)} ativa a tabela de caracteres em que o
    código 248 é o símbolo de grau ($^\circ$).
  \item Texto branco, tamanho 2 (cada caractere com 12$\times$16 pixels).
\end{itemize}

\subsection{\texttt{loop}: executado sem parar}
\begin{lstlisting}
void loop() {
  float ref = le(PINO_REF);
  float novo = atan2(leG(PINO_Y, ref), leG(PINO_Z, ref)) * RAD_TO_DEG;
  angulo += ALFA * (novo - angulo);
  if (fabs(angulo - mostrado) > 0.8) mostrado = round(angulo);

  digitalWrite(LED, abs(mostrado) >= LIMIAR);

  tela.clearDisplay();
  tela.setCursor(0, 24);
  tela.print(F("X  "));
  tela.print(mostrado);
  tela.write(248);
  tela.display();

  delay(50);
}
\end{lstlisting}
Linha a linha:
\begin{enumerate}
  \item Lê a referência de 3,3~V uma vez por ciclo.
  \item Lê Y e Z em g e calcula o ângulo. \texttt{atan2} devolve
    radianos; \texttt{RAD\_TO\_DEG} ($180/\pi$) converte para graus.
  \item Filtro passa-baixa (Seção~\ref{sec:filtro}).
  \item Histerese (Seção~\ref{sec:hist}). \texttt{fabs} é o valor
    absoluto para \texttt{float}; \texttt{round} arredonda para o
    inteiro mais próximo.
  \item O LED recebe o resultado da comparação: verdadeiro (1, acende)
    se o módulo do ângulo for pelo menos 30, falso (0, apaga) se não.
    \texttt{abs} cobre os dois lados da inclinação.
  \item O desenho é feito numa memória interna (\emph{buffer}): limpa,
    posiciona o cursor no meio da altura (linha 24), escreve
    ``X'', o número e o símbolo de grau. \texttt{F(\dots)} guarda o
    texto fixo na memória de programa, poupando a RAM de 2~kB do Uno.
  \item Só \texttt{tela.display()} envia o buffer pelo
    I\textsuperscript{2}C. Assim a tela troca de uma vez e não pisca.
  \item Espera 50~ms antes do próximo ciclo.
\end{enumerate}

\section{Cuidados de hardware}\label{sec:hw}
O OLED e o GY-61 estão em 3,3~V, mas a \texttt{Wire} liga resistores
internos (\emph{pull-ups}) que puxam SDA e SCL para os 5~V do Uno. Isso
põe 5~V nos pinos de um dispositivo de 3,3~V. O programa os desliga
logo depois de iniciar a tela. Com isso, o barramento depende dos
pull-ups do próprio módulo OLED, ligados a 3,3~V; o ATmega328P lê
3,3~V como nível alto, então a comunicação funciona.

\section{Calibração}\label{sec:calib}
512 e 102,4 são os valores teóricos. Cada ADXL335 tem pequenas
diferenças de zero e de sensibilidade. Para calibrar:
\begin{enumerate}
  \item Imprimir no monitor serial o valor corrigido
    $N/N_\text{A3}\times 1024$ de Y e Z.
  \item Medir cada eixo apontado para cima ($+1$~g) e para baixo
    ($-1$~g).
  \item \texttt{ZERO} é a média das duas leituras; \texttt{ESCALA} é a
    metade da diferença entre elas.
\end{enumerate}
O código usa um só par de constantes para os dois eixos; se Y e Z
diferirem muito, cada eixo precisaria das suas.

\section{Limitações}
\begin{itemize}
  \item O acelerômetro não distingue gravidade de movimento: o ângulo
    só é correto com o sensor parado. Vibrar ou chacoalhar gera erro.
  \item O filtro que estabiliza a leitura também a atrasa
    (Seção~\ref{sec:filtro}).
  \item Com o sensor em pé (eixo X na vertical), $a_y$ e $a_z$ ficam
    perto de zero e o ângulo perde o sentido.
  \item Perto de $\pm180^\circ$ o \texttt{atan2} salta de $+180$ para
    $-180$, e o filtro passaria pelos valores intermediários. Na faixa
    de uso (bem abaixo disso) não ocorre.
\end{itemize}

\section{Conclusão}
O inclinômetro atende aos critérios da Tabela~\ref{tab:criterios}: mede
o ângulo em X sem depender da tensão do USB, graças à leitura dos
3,3~V em A3; mostra um valor estável graças à média, ao filtro e à
histerese; e acende o LED a 30$^\circ$ para os dois lados. O preço da
estabilidade é um atraso de cerca de um a dois segundos na resposta.

\end{document}
